% year: תשפ"ד
% type: מבחן
% semester: א
% moed: ב
% number: 3
% points: 14
% categories: func-analysis, derivatives
% source: exams/מבחנים/תשפד/מועד ב תשפד.pdf

\begin{question}
חשבו תחומי קמירות כלפי מעלה וקמירות כלפי מטה (קמירות וקעירות), ונקודות פיתול של הפונקציה
\[ f(x)=\ln\left(\frac{e^x}{1+e^x}\right) \]
בתחום הגדרתה.
\end{question}

\begin{hint}
פשטו קודם: $\ln\frac{e^x}{1+e^x}=x-\ln(1+e^x)$.
\end{hint}

\begin{solution}
\textbf{תחום הגדרה.} לכל $x$ ממשי $\frac{e^x}{1+e^x}>0$, ולכן $f$ מוגדרת על כל $\R$. לפי חוקי לוגריתמים
\[ f(x)=\ln e^x-\ln(1+e^x)=x-\ln(1+e^x). \]

\textbf{נגזרת ראשונה.}
\[ f'(x)=1-\frac{e^x}{1+e^x}=\frac{1}{1+e^x}. \]

\textbf{נגזרת שנייה.} לפי כלל השרשרת,
\[ f''(x)=-\frac{e^x}{(1+e^x)^2}. \]
לכל $x\in\R$ מתקיים $e^x>0$ ו-$(1+e^x)^2>0$, ולכן $f''(x)<0$ לכל $x$.

\textbf{מסקנה.} מאחר ש-$f''<0$ על כל $\R$, הפונקציה קמורה כלפי מטה (קעורה, $\cap$) על כל הישר $(-\infty,\infty)$, ואין תחום שבו היא קמורה כלפי מעלה. נקודת פיתול היא נקודה שבה משתנה כיוון הקמירות; מאחר ש-$f''$ אינה מחליפה סימן (ואף אינה מתאפסת), \textbf{אין לפונקציה נקודות פיתול}.
\end{solution}
